\chapter*{General Introduction}
\label{ch:intoGen}

Deep reinforcement learning (DRL) has been applied to sequential decision-making tasks, including games and robotic control. However, standard loss functions remain sensitive to the scale of rewards encountered during training. When reward magnitudes vary across environments or change dynamically, regression losses such as mean squared error (MSE) can produce large gradient updates, destabilizing the training process and causing value function divergence.

This challenge is significant in \emph{imagination-based} reinforcement learning. Model-based approaches, such as the Dreamer family of agents, can improve sample efficiency by allowing training on simulated experience. However, because imagined rollouts are generated autoregressively, small prediction errors compound across time steps. Combined with scale-sensitive losses, this compounding effect can cause predicted values and rewards to diverge, leading to a value explosion.

DreamerV3 addressed this problem by introducing \emph{symlog}-transformed losses. The symlog function, defined as $\operatorname{symlog}(x) = \operatorname{sign}(x) \cdot \ln(|x| + 1)$, compresses large magnitudes logarithmically while preserving the sign. By applying this transformation to critic and reward prediction targets before computing the MSE loss, DreamerV3 achieved scale-invariant learning without requiring task-specific hyperparameter tuning.

This project investigates whether this symlog transformation retains its stabilizing properties when integrated into a simplified imagination-based agent. Specifically, we implement a simplified Dreamer-like agent in PyTorch that operates in the raw state space rather than a learned latent space. The agent uses a Dyna-style training scheme, learning from both real environment interactions and imagined trajectories generated by a learned world model. The actor, critic, and world model components are implemented as two-layer multilayer perceptrons with 256 hidden units each.

We compare three configurations:
\begin{itemize}
    \item \textbf{Condition A (Model-Free Actor-Critic):} A baseline agent that learns exclusively from real environment transitions.
    \item \textbf{Condition B (Imagination AC with Standard MSE):} An imagination-augmented agent that trains on both real and imagined data, using standard MSE losses.
    \item \textbf{Condition C (Imagination AC with Symlog MSE):} An imagination-augmented agent using symlog-transformed targets in the critic and reward prediction losses.
\end{itemize}

\noindent We evaluate these configurations on CartPole-v1 and LunarLander-v3 using five random seeds, yielding 30 runs.

The results reveal distinct differences between the configurations. Without symlog transformations, imagination-augmented training leads to value divergence: critic losses exceed $10^{15}$, causing policy collapse. Symlog-transformed losses prevent this divergence, maintaining stable losses throughout training. These findings indicate that symlog is a broadly applicable technique for stabilizing imagination-based reinforcement learning.


%%%%%%%%%%%%%%%%%%%%%%%%%%%%%%%%%%%%%%%%%%%%%%%%%%%%%%%%%%%%%%%%%%%%%%%%%%%%%%%%%%%
\section*{Organisation of the Report}
\label{sec:intro_org}

This report contains seven chapters. Chapter~\ref{ch:into} introduces the background, problem statement, objectives, and solution approach. Chapter~\ref{ch:lit_rev} reviews literature on model-free and model-based reinforcement learning, the Dreamer family, and the symlog transformation. Chapter~\ref{ch:method} details the methodology, including agent architectures, training procedures, and environments. Chapter~\ref{ch:results} presents the experimental results. Chapter~\ref{ch:evaluation} discusses the findings, study limitations, and broader context. Chapter~\ref{ch:con} summarizes the conclusions and proposes future work. Finally, Chapter~\ref{ch:reflection} provides a reflection on the project.
